\documentclass[10pt,letterpaper]{article}
\usepackage[utf8]{inputenc}
\usepackage[spanish]{babel}
\usepackage[T1]{fontenc}
\usepackage{geometry}
\usepackage{enumitem}
\usepackage{titlesec}
\usepackage[hidelinks]{hyperref}

\geometry{left=0.45in, right=0.45in, top=0.42in, bottom=0.42in}

\titleformat{\section}{\small\bfseries\uppercase}{}{0pt}{}[\titlerule]
\titlespacing{\section}{0pt}{6pt}{3pt}

\setlist[itemize]{topsep=0pt, parsep=0pt, partopsep=0pt, itemsep=0pt, leftmargin=1.35em}
\linespread{0.95}
\setlength{\parindent}{0pt}

\begin{document}
\pagestyle{empty}

\begin{center}
	\textbf{\Huge DIEGO VÉLIZ} \\ \vspace{4pt}
	Buin, Región Metropolitana \ | \ +56 9 3190 7792 \ |
	\ \href{mailto:diego.veliz.valderrama@gmail.com}
	{diego.veliz.valderrama@gmail.com} \\
	\href{https://linkedin.com/in/diego-véliz-valderrama}{linkedin.com/in/diego-véliz-valderrama}
	\ | \ \href{https://github.com/eldiegoveliz}{github.com/eldiegoveliz} \
	| \ \href{https://diegoveliz.xyz}{diegoveliz.xyz}
\end{center}

\section{Perfil}
Licenciado en Economía y estudiante de Ingeniería Comercial, enfocado en
\textbf{equity research}, valoración y análisis financiero. Desarrollo modelos
de valoración, tesis de inversión y herramientas de automatización en Python
y Excel.

\section{Educación}
\textbf{Universidad de los Andes} \hfill \textbf{Las Condes, Chile} \\
\textit{Ingeniería Comercial, mención en Economía} \hfill
\textit{Esperado ago. 2027} \\
\textit{Licenciatura en Economía} \hfill
\textit{Obtenida ago. 2026}
\begin{itemize}
	\item \textbf{Beca Puntaje PTU}.
	\item \textbf{Minor:} Neurociencias.
	\item \textbf{Cursos relevantes:} Renta Fija, Economía Fiscal,
		Investment Management and Research, Advanced Corporate Finance.
\end{itemize}

\vspace{4pt}
\textbf{Colegio Técnico Profesional Nocedal}
\hfill \textbf{Santiago, Chile} \\
\textit{Técnico Medio en Electricidad y Electrónica}
\hfill \textit{2007 -- 2019}

\section{Experiencia académica y liderazgo}
\textbf{Universidad de los Andes} \hfill \textbf{Santiago, Chile} \\
\textit{Ayudante de cátedra} \hfill \textit{jul. 2022 -- dic. 2025}
\begin{itemize}
	\item \textbf{Microeconomía II:} Diseñé materiales didácticos
		y pautas autoexplicativas de resolución de problemas
		(teoría neoclásica de producción, internalización y contrato)
		para cohortes de \textbf{más de 100 estudiantes}.
	\item \textbf{Álgebra Lineal y Excel I:} Dicté ayudantías técnicas, evalué
		tareas/pruebas y reforcé herramientas cuantitativas para cursos de
		\textbf{más de 130 estudiantes}.
	\item \textbf{Productividad e Incentivos:} Analicé tareas semanales basadas
		en literatura académica y modelos de incentivos para
		\textbf{más de 100 estudiantes}.
\end{itemize}

\vspace{2pt}
\textbf{Club de Economía UANDES} \hfill \textbf{Santiago, Chile} \\
\textit{Cofundador} \hfill \textit{ago. 2024 -- ago. 2026}
\begin{itemize}
	\item Cofundé una organización estudiantil para acercar teoría económica,
		mercados y práctica profesional mediante charlas, visitas y redes con
		la industria.
\end{itemize}

\section{Research financiero y proyectos}
\textbf{Modelos de valoración y tesis de inversión}
\begin{itemize}
	\item Elaboré valoraciones y escenarios de inversión para
		\textbf{Abivax, Sarepta, Spruce, LMT y Beeline}; publiqué modelos
		y notas de análisis en
		\underline{\href{https://diegoveliz.xyz}{diegoveliz.xyz}}.
	\item Analicé oportunidades de inversión en biotech considerando evidencia
		clínica, mercado potencial, riesgos regulatorios y perspectivas
		comerciales.
	\item Proyecté escenarios clínicos de Obefazimod y su impacto en la
		valoración de \textbf{Abivax}.
\end{itemize}

\vspace{2pt}
\textbf{Simulador long/short y análisis de empresas biotech}
\begin{itemize}
	\item Desarrollé un simulador para estudiar estrategias long/short,
		construcción de carteras y desempeño fuera de muestra.
	\item Incorporé un universo de \textbf{582 empresas biotecnológicas},
		agrupadas por indicación terapéutica, para identificar compañías de
		interés según su fase clínica y resultados históricos.
	\item Publiqué una aplicación web con actualización semanal y acceso al
		código fuente:
		\href{https://simulador.diegoveliz.xyz/biotech/}{Aplicación web} |
		\href{https://github.com/eldiegoveliz/biotech-hf-calculator}{GitHub}.
\end{itemize}

\vspace{2pt}
\textbf{Monitoreo de mercados y catalizadores biotech}
\begin{itemize}
	\item Desarrollé un monitor de noticias y eventos clínicos y regulatorios
		para seguir oportunidades en empresas biotech de pequeña capitalización.
	\item Incorporé seguimiento de precios y cotizaciones de compra y venta en
		tiempo real para acciones seleccionadas.
\end{itemize}

\section{Certificaciones y formación complementaria}
\begin{itemize}
	\item \textbf{AI Leadership}, OpenAI Academy \hfill \textit{sept. 2026}
	\item \textbf{Python}, Kaggle \hfill \textit{sept. 2026}
	\item \textbf{Programa Andino}, ESE Business School \hfill \textit{nov. 2024}
	\item \textbf{Gestión de organizaciones efectivas}, PUC de Chile/Coursera \hfill
		\textit{sept. 2020}
\end{itemize}

\section{Habilidades e intereses}
\begin{itemize}
	\item \textbf{Idiomas:} Inglés (fluido), español (nativo).
	\item \textbf{Herramientas:} Excel (avanzado), Bloomberg Terminal,
		Python (básico), Git, Linux/Nginx, SQL/SQLite (básico) y \LaTeX.
	\item \textbf{Intereses:} Equity research, renta fija, biotech,
		macroeconomía y automatización de research.
\end{itemize}

\end{document}
